% ============================================================
% Savoury -- Test Cases
% Follows the "Test Case Template (Doc:T_01)" layout:
% header box, pre-conditions, steps table, post-conditions.
% Compile with pdfLaTeX. Compile twice so the "Page: 1 of N"
% field is filled in (Overleaf does this automatically).
% ============================================================
\documentclass[11pt]{article}

\usepackage[a4paper,landscape,margin=1.8cm]{geometry}
\usepackage[T1]{fontenc}
\usepackage[utf8]{inputenc}
\usepackage{array}
\usepackage{tabularx}
\usepackage{longtable}
\usepackage[table]{xcolor}
\usepackage{refcount}
\usepackage{fancyhdr}

\setlength{\parindent}{0pt}
\renewcommand{\arraystretch}{1.25}

\newcolumntype{L}[1]{>{\raggedright\arraybackslash}p{#1}}
\newcolumntype{C}[1]{>{\centering\arraybackslash}p{#1}}
\definecolor{headgray}{gray}{0.88}

% Running header, as in the template
\pagestyle{fancy}
\fancyhf{}
\lhead{\small Savoury -- Software Engineering (UCS503)}
\rhead{\small Page \thepage}
\renewcommand{\headrulewidth}{0pt}

% ---------- Template building blocks ----------
\newcommand{\val}[1]{\textcolor{blue}{#1}}
\newcommand{\tcpages}[1]{\number\numexpr\getpagerefnumber{tc:#1:end}-\getpagerefnumber{tc:#1:start}+1\relax}

% \tcheader{number}{name}{subsystem}{short description}
\newcommand{\tcheader}[4]{%
  \clearpage
  \def\tcid{#1}%
  \noindent\label{tc:#1:start}%
  \begin{tabularx}{\linewidth}{|L{0.36\linewidth}X L{3.4cm}|}
    \hline
    \textbf{Test Case \#:} \val{#1} &
    \textbf{Test Case Name:} \val{#2} &
    \textbf{Page:} \val{\textbf{1 of \tcpages{#1}}} \\
    \textbf{System:} \val{\textbf{Savoury}} &
    \textbf{Subsystem:} \val{\textbf{#3}} & \\
    \textbf{Designed by:} \val{\textbf{Savoury Team}} &
    \textbf{Design Date:} \val{\textbf{30/10/2026}} & \\
    \textbf{Executed by:} &
    \textbf{Execution Date:} & \\
    \multicolumn{3}{|L{\dimexpr\linewidth-2\tabcolsep-2\arrayrulewidth\relax}|}%
      {\textbf{Short Description:} \val{\textbf{#4}}} \\
    \hline
  \end{tabularx}\par\medskip}

% \tcbox{title}{content}  -- pre/post-conditions box
\newcommand{\tcbox}[2]{%
  \noindent\begin{tabularx}{\linewidth}{|X|}
    \hline
    \textbf{#1}\newline #2 \\
    \hline
  \end{tabularx}\par\medskip}

% Steps table
\newcounter{tcstep}
\newcommand{\tcstep}[2]{\stepcounter{tcstep}\thetcstep & #1 & #2 & & \\ \hline}
\newcommand{\tccheck}[1]{\stepcounter{tcstep}\thetcstep & \textbf{Check post-condition #1} & & & \\ \hline}
\newenvironment{tcsteps}{%
  \setcounter{tcstep}{0}\small
  \begin{longtable}{|C{0.8cm}|L{7.0cm}|L{10.4cm}|C{1.1cm}|L{3.9cm}|}
    \hline
    \rowcolor{headgray}\textbf{Step} & \textbf{Action} &
      \textbf{Expected System Response} & \textbf{Pass/ Fail} & \textbf{Comment} \\
    \hline
    \endhead}
  {\end{longtable}}

% Post-conditions box that also closes the test case
\newcommand{\tcend}[1]{\tcbox{Post-conditions}{#1}\label{tc:\tcid:end}}

\begin{document}

% ============================================================
% Title page
% ============================================================
\thispagestyle{empty}
\begin{center}
  \vspace*{2.5cm}
  {\fontsize{40}{48}\selectfont\bfseries Savoury\par}
  \vspace{0.6cm}
  {\fontsize{30}{36}\selectfont\bfseries Test Cases\par}
  \vspace{0.8cm}
  {\Large Hostel Mess Management System\par}
  \vspace{0.4cm}
  {\large Based on the Test Case Template (Doc:T\_01)\par}
  \vspace{1.5cm}
  {\large Siddharth Singla (1024170252) \quad
   Vineet Tripathi (1024170253) \quad
   Shivansh Sahu (1024170249)\par}
  \vspace{0.5cm}
  {\large UCS503 Software Engineering, Thapar Institute of Engineering \& Technology, Patiala\par}
\end{center}

% ============================================================
% Test environment and data
% ============================================================
\clearpage
\section*{Test Environment and Test Data}

All test cases use the accounts and settings below. \textbf{D} is the date the
test is run. All times are Indian Standard Time (IST). The password for every
test account is \texttt{Test@1234}. Each test case lists its own
pre-conditions, so test cases can be run on different days; set up the
pre-conditions before each one.

\medskip
\textbf{Clients and tools:} web console in a desktop browser; the Savoury
mobile app on an Android phone set to IST; Postman for steps that call the API
directly. API steps send the header \texttt{Authorization: Bearer <token>}, where
the token comes from \texttt{POST /auth/login}. API dates are sent as
\texttt{<date>T12:00:00Z}.

\medskip
\begin{tabularx}{\linewidth}{|L{4.2cm}|L{5.4cm}|L{4.2cm}|L{2.4cm}|X|}
  \hline
  \rowcolor{headgray}\textbf{Name} & \textbf{Email} & \textbf{Role} & \textbf{Hostel} & \textbf{Notes} \\ \hline
  Super Admin  & admin@thapar.edu       & Super Admin    & --       & Created on the Setup page \\ \hline
  Neha Gupta   & warden.a@thapar.edu    & Warden         & Hostel A & Created in TC 2.1 \\ \hline
  Arjun Mehta  & arjun@thapar.edu       & Co-Warden      & Hostel A & Created in TC 2.1 \\ \hline
  Ravi Singh   & caretaker.a@thapar.edu & Caretaker      & Hostel A & Created in TC 2.1 \\ \hline
  Suresh Yadav & staff.a@thapar.edu     & Mess Staff     & Hostel A & Created in TC 2.2 \\ \hline
  Rahul Kumar  & rahul@thapar.edu       & Student        & Hostel A & Roll 102103001, wallet Rs.~1000.00 \\ \hline
  Priya Sharma & priya@thapar.edu       & Student        & Hostel A & Roll 102103002, wallet Rs.~1000.00 \\ \hline
  Karan Patel  & karan@thapar.edu       & Student, then Mess Committee & Hostel A & Roll 102103003, wallet Rs.~1000.00; promoted in TC 2.2 \\ \hline
  Dev Malhotra & dev@thapar.edu         & Student        & Hostel A & Roll 102103004, no wallet until TC 6.1 \\ \hline
  Aman Verma   & aman@thapar.edu        & Student        & Hostel B & Roll 102203001, wallet Rs.~1000.00 \\ \hline
\end{tabularx}

\medskip
Wallets are for the semester \texttt{Autumn Semester 2026}.

\medskip
\begin{tabularx}{\linewidth}{|L{4cm}|L{5cm}|L{5cm}|X|}
  \hline
  \rowcolor{headgray}\textbf{Meal} & \textbf{Serving window} & \textbf{Cutoff} & \textbf{Rate} \\ \hline
  Breakfast & 07:00 -- 09:30 & 00:00 on the same day & Rs.~75.00 \\ \hline
  Lunch     & 12:00 -- 14:30 & 09:00 on the same day & Rs.~75.00 \\ \hline
  Dinner    & 19:00 -- 21:30 & 14:00 on the same day & Rs.~75.00 \\ \hline
\end{tabularx}

\medskip
Counter check-in is accepted from 60 minutes before a meal's serving window
opens until 60 minutes after it closes, for today's meals only.

% ============================================================
% Index
% ============================================================
\clearpage
\section*{Index of Test Cases}

\begin{tabularx}{\linewidth}{|L{1.6cm}|L{6.5cm}|L{4.2cm}|L{3cm}|X|}
  \hline
  \rowcolor{headgray}\textbf{\#} & \textbf{Test Case Name} & \textbf{Subsystem} & \textbf{Client} & \textbf{Requirements (SRS 3.2)} \\ \hline
  1.1 & Student Registration                 & Authentication  & Mobile        & FR-2 \\ \hline
  1.2 & Staff Portal Login                   & Authentication  & Web           & FR-3, FR-4 \\ \hline
  2.1 & Appoint Hostel Administrators        & User Management & Web           & FR-8, FR-10 \\ \hline
  2.2 & Add Mess Staff and Role Locks        & User Management & Web           & FR-9, FR-10, FR-11 \\ \hline
  3.1 & Opt Out of a Single Meal             & Meal Booking    & Mobile, API   & FR-14, FR-15, FR-17 \\ \hline
  3.2 & Opt Out Across a Date Range          & Meal Booking    & Mobile        & FR-18, FR-19 \\ \hline
  4.1 & Serve a Student at the Counter       & Counter Check-In & Web, Mobile  & FR-21, FR-22 \\ \hline
  4.2 & Check-In Rules Enforced by the Server & Counter Check-In & API         & FR-4, FR-23, FR-24, FR-26, FR-27 \\ \hline
  4.3 & Walk-in Override                     & Counter Check-In & API          & FR-25 \\ \hline
  5.1 & Live Headcount Dashboard             & Headcount       & Web, Mobile   & FR-28, FR-29, FR-30 \\ \hline
  6.1 & Wallet Top-Up                        & Wallet          & Web, Mobile   & FR-31, FR-32 \\ \hline
  6.2 & Semester-End Cashout                 & Wallet          & Mobile, Web, API & FR-33, FR-34, FR-35 \\ \hline
  7.1 & Hostel-Scoped Notices                & Notices         & Web, Mobile   & FR-39, FR-40, FR-41 \\ \hline
\end{tabularx}

% ============================================================
% 1.1 Student Registration
% ============================================================
\tcheader{1.1}{Student Registration}{Authentication}
  {Test that a student can register on the mobile app with a college email (FR-2)}

\tcbox{Pre-conditions}{%
  The mobile app is open on the Sign In screen.\newline
  Hostel B exists.\newline
  No account exists for isha@thapar.edu.}

\begin{tcsteps}
\tcstep{Open the registration screen from the Sign In screen}
       {The system displays the registration form with name, email, password, roll number, room number, phone and a `Select Hostel' list}
\tcstep{Enter name `Isha Arora', email `isha@gmail.com', password `Test@1234', roll `102103010', room `B-214', phone `9876543210', select `Hostel B' and tap `Create Account'}
       {The system shows a `Registration Failed' message saying a college (@thapar.edu) email is required. No account is created}
\tcstep{Change the email to `isha@thapar.edu' and the phone to `98765', then tap `Create Account'}
       {The system shows a `Registration Failed' message saying the phone number must be exactly 10 digits}
\tcstep{Change the phone to `9876543210' and tap `Create Account'}
       {The account is created, the student is signed in, and the home screen with the menu is displayed}
\tccheck{1}
\tcstep{Sign out, open the registration screen again and repeat step 4 with the same details}
       {The system shows a `Registration Failed' message saying the email is already in use}
\tccheck{2}
\end{tcsteps}

\tcend{%
  1. A user `Isha Arora' with email isha@thapar.edu, role Student, Hostel B and roll 102103010 is saved in the database.\newline
  2. Only one account exists with the email isha@thapar.edu.}

% ============================================================
% 1.2 Staff Portal Login
% ============================================================
\tcheader{1.2}{Staff Portal Login}{Authentication}
  {Test login to the web console, including wrong input and student accounts (FR-3, FR-4)}

\tcbox{Pre-conditions}{%
  The web console is open on the login page.\newline
  The Warden account warden.a@thapar.edu (Hostel A) and the Student account rahul@thapar.edu exist, both with password `Test@1234'.}

\begin{tcsteps}
\tcstep{Leave both fields empty and click `Sign In'}
       {The system displays `Both fields are required.'}
\tcstep{Enter `warden.a@gmail.com' and `Test@1234', click `Sign In'}
       {The system displays `Email must end with @thapar.edu.'}
\tcstep{Enter `warden.a@thapar.edu' and `wrongpass', click `Sign In'}
       {The system displays `Invalid email or password.'}
\tcstep{Enter `rahul@thapar.edu' and `Test@1234', click `Sign In'}
       {The system displays `Students cannot access the staff portal. Please use the mobile app.' and stays on the login page}
\tcstep{Enter `warden.a@thapar.edu' and `Test@1234', click `Sign In'}
       {The system opens the Check-In page. The sidebar shows Check-In, Headcount, Menu, Users, Wallet and Notices, but not Hostels}
\tccheck{1}
\tcstep{Click the sign-out icon next to the user's name, then open the \texttt{/users} address directly}
       {The system shows the login page in both cases}
\tccheck{2}
\end{tcsteps}

\tcend{%
  1. The browser holds a session token for Neha Gupta with role Warden and Hostel A.\newline
  2. After sign-out, the browser holds no session token.}

% ============================================================
% 2.1 Appoint Hostel Administrators
% ============================================================
\tcheader{2.1}{Appoint Hostel Administrators}{User Management}
  {Test that the Super Admin can appoint one Warden, one Co-Warden and Caretakers per hostel (FR-8, FR-10)}

\tcbox{Pre-conditions}{%
  The user is logged in to the web console as the Super Admin and the Users page is open.\newline
  Hostel A exists and has no Warden, Co-Warden or Caretaker.}

\begin{tcsteps}
\tcstep{Click `Add User'}
       {The system displays the Add User form with name, email, password, role, hostel and phone fields}
\tcstep{Enter `Neha Gupta', `warden.a@thapar.edu', `Test@1234', choose role `Warden', leave hostel unselected and click `Create User'}
       {The system displays `Please select a hostel for this user' and does not create the account}
\tcstep{Select hostel `Hostel A' and click `Create User'}
       {The form closes and Neha Gupta appears in the list with a `Warden' badge and Hostel A}
\tccheck{1}
\tcstep{Click `Add User', enter `Arjun Mehta', `arjun@thapar.edu', `Test@1234', role `Warden', hostel `Hostel A', click `Create User'}
       {The form stays open and displays `This hostel already has a Warden. Remove the existing one first.'}
\tcstep{Change the role to `Co-Warden' and click `Create User'}
       {The form closes and Arjun Mehta appears with a `Co-Warden' badge}
\tcstep{Click `Add User', enter `Ravi Singh', `caretaker.a@thapar.edu', `Test@1234', role `Caretaker', hostel `Hostel A', click `Create User'}
       {The form closes and Ravi Singh appears with a `Caretaker' badge}
\tccheck{2}
\end{tcsteps}

\tcend{%
  1. Neha Gupta is saved with role Warden in Hostel A.\newline
  2. Hostel A has exactly one Warden (Neha Gupta), one Co-Warden (Arjun Mehta) and one Caretaker (Ravi Singh).}

% ============================================================
% 2.2 Add Mess Staff and Role Locks
% ============================================================
\tcheader{2.2}{Add Mess Staff and Role Locks}{User Management}
  {Test that a Warden can add Mess Staff and switch students to Mess Committee, and that a Caretaker cannot add staff (FR-9, FR-10, FR-11)}

\tcbox{Pre-conditions}{%
  Hostel A has the Warden Neha Gupta, the Caretaker Ravi Singh and the students Priya Sharma and Karan Patel (role Student).\newline
  No account exists for staff.a@thapar.edu.\newline
  The user is logged in to the web console as warden.a@thapar.edu.}

\begin{tcsteps}
\tcstep{Open the Users page}
       {The system lists only Hostel A users and shows an `Add Mess Staff' button}
\tcstep{Click `Add Mess Staff', enter `Suresh Yadav', `staff.a@thapar.edu', `Test@1234' and click `Create User'}
       {The role field shows `Mess Staff' and cannot be changed. The form closes and Suresh Yadav appears with a `Mess Staff' badge}
\tccheck{1}
\tcstep{Look at the Role \& Access column for Suresh Yadav, Ravi Singh and Neha Gupta}
       {Each role is shown as a fixed badge with no drop-down, so it cannot be changed}
\tcstep{Change Karan Patel's role drop-down from `Student' to `Mess Sec.'}
       {Karan Patel's badge changes to `Mess Sec.'}
\tcstep{Change Priya Sharma's role drop-down from `Student' to `Mess Staff'}
       {The system displays `You can only switch student accounts between Student and Mess Committee.' and Priya stays a Student}
\tccheck{2}
\tcstep{Sign out and log in as caretaker.a@thapar.edu, then open the Users page}
       {The system lists Hostel A users and shows no `Add Mess Staff' button}
\end{tcsteps}

\tcend{%
  1. Suresh Yadav is saved with role Mess Staff in Hostel A.\newline
  2. Karan Patel's role is Mess Committee and Priya Sharma's role is still Student.}

% ============================================================
% 3.1 Opt Out of a Single Meal
% ============================================================
\tcheader{3.1}{Opt Out of a Single Meal}{Meal Booking}
  {Test opting out of and back into one meal before its cutoff, and that the meal is closed after the cutoff (FR-14, FR-15, FR-17)}

\tcbox{Pre-conditions}{%
  Rahul Kumar is signed in to the mobile app.\newline
  The current time on day D is between 07:00 and 08:59, before the 09:00 Lunch cutoff.\newline
  A menu is published for Lunch on day D, and Rahul has not changed his Lunch booking for day D.}

\begin{tcsteps}
\tcstep{Open the home screen and look at the Lunch card for day D}
       {The card shows `You're opted in' and `Tap to opt out'}
\tcstep{Tap the Lunch toggle}
       {The card shows `Opted out' and `Tap to opt back in'}
\tccheck{1}
\tcstep{Tap the Lunch toggle again}
       {The card shows `You're opted in'}
\tcstep{Tap the Lunch toggle once more}
       {The card shows `Opted out'}
\tccheck{2}
\tcstep{After 09:00, reopen the app and look at the Lunch card}
       {The card shows `Opted out' and `Closed $\cdot$ You will not be served', and there is no toggle}
\tcstep{In Postman, as Rahul, send \texttt{PUT /bookings/toggle} for Lunch on day D with status \texttt{OPTED\_IN}}
       {The system returns HTTP 400 with `Cutoff time has passed for this meal'}
\tccheck{3}
\end{tcsteps}

\tcend{%
  1. Rahul's Lunch booking for day D is saved with status \texttt{OPTED\_OUT}.\newline
  2. Exactly one booking exists for Rahul, Lunch and day D, with status \texttt{OPTED\_OUT}.\newline
  3. That booking is still \texttt{OPTED\_OUT} after the cutoff.}

% ============================================================
% 3.2 Opt Out Across a Date Range
% ============================================================
\tcheader{3.2}{Opt Out Across a Date Range}{Meal Booking}
  {Test scheduling an absence for several days, cancelling one day, and skipping meals past their cutoff (FR-18, FR-19)}

\tcbox{Pre-conditions}{%
  Rahul Kumar is signed in to the mobile app.\newline
  The current time on day D is between 09:00 and 13:59, so Breakfast and Lunch on day D are past their cutoff and Dinner is not.\newline
  Rahul has no scheduled absences from day D to day D+4.}

\begin{tcsteps}
\tcstep{Open the Opt-Out Calendar}
       {The system displays the meal absence planner and `No Upcoming Absences'}
\tcstep{Select day D+2, then day D+4}
       {Days D+2 to D+4 are highlighted and the button reads `Opt Out (3d)'}
\tcstep{Tap `Opt Out (3d)', then tap `Confirm'}
       {The system displays `Opt-Out Confirmed' with `Successfully opted out of 9 open meal(s) across 3 day(s).' The three days are listed under scheduled absences}
\tccheck{1}
\tcstep{Tap `Re-Opt In' on day D+3 and tap `Confirm'}
       {The system displays `Re-opted In' and day D+3 is removed from the list}
\tccheck{2}
\tcstep{Select only day D, tap `Opt Out (1d)', then tap `Confirm'}
       {The system displays `Opt-Out Confirmed' with `Successfully opted out of 1 open meal(s) across 1 day(s).' and a note that 2 meal(s) whose cutoff had passed were kept as-is}
\tccheck{3}
\end{tcsteps}

\tcend{%
  1. Absences are saved for days D+2, D+3 and D+4, and all 9 of their meals are \texttt{OPTED\_OUT}.\newline
  2. The absence for day D+3 is removed and its 3 meals are \texttt{OPTED\_IN} again. Days D+2 and D+4 are unchanged.\newline
  3. Dinner on day D is \texttt{OPTED\_OUT}. Breakfast and Lunch on day D are unchanged.}

% ============================================================
% 4.1 Serve a Student at the Counter
% ============================================================
\tcheader{4.1}{Serve a Student at the Counter}{Counter Check-In}
  {Test that serving a student records the meal and deducts the meal rate from the wallet (FR-21, FR-22)}

\tcbox{Pre-conditions}{%
  The user is logged in to the web console as Mess Staff staff.a@thapar.edu.\newline
  The current time on day D is between 12:00 and 14:30 (Lunch is being served).\newline
  Rahul Kumar is opted in to Lunch on day D and has a wallet balance of Rs.~1000.00.\newline
  Priya Sharma has opted out of Lunch on day D.}

\begin{tcsteps}
\tcstep{Open the Check-In page with Lunch selected}
       {The system shows `Service Active'. Rahul Kumar is listed with a `Tap to Serve' button. Priya Sharma is not listed}
\tcstep{Type `102103001' in `Search records...'}
       {Only Rahul Kumar is listed}
\tcstep{Click `Tap to Serve' for Rahul Kumar}
       {Rahul leaves the waiting list and appears in the recent check-ins with a timestamp. The served count increases by 1}
\tcstep{Wait 5 seconds and reload the page}
       {Rahul Kumar is still not in the waiting list}
\tccheck{1}
\tcstep{Sign in to the mobile app as Rahul Kumar and open Wallet}
       {The balance shows Rs.~925.00 and the transaction history shows a `Meal Deduction' of Rs.~75.00}
\tccheck{2}
\end{tcsteps}

\tcend{%
  1. An attendance record is saved for Rahul's Lunch on day D with result \texttt{SERVED}, checked in by Suresh Yadav.\newline
  2. Rahul's wallet balance is Rs.~925.00 and a \texttt{MEAL\_DEDUCTION} transaction of Rs.~75.00 is saved for that attendance.}

% ============================================================
% 4.2 Check-In Rules Enforced by the Server
% ============================================================
\tcheader{4.2}{Check-In Rules Enforced by the Server}{Counter Check-In}
  {Test that the API rejects duplicate, opted-out, other-hostel, no-wallet and out-of-hours check-ins (FR-4, FR-23, FR-24, FR-26, FR-27)}

\tcbox{Pre-conditions}{%
  Postman holds a token for Mess Staff staff.a@thapar.edu (Hostel A) and the ids of the Lunch and Dinner meal types.\newline
  The current time on day D is between 12:00 and 14:30.\newline
  Rahul Kumar has already been served Lunch on day D (Test Case 4.1) and his balance is Rs.~925.00.\newline
  Priya Sharma has opted out of Lunch on day D. Aman Verma is in Hostel B. Dev Malhotra (Hostel A) has no wallet.}

\begin{tcsteps}
\tcstep{Send \texttt{POST /attendance/mark-served} for Rahul Kumar, Lunch, day D}
       {HTTP 200 with attendance result \texttt{DUPLICATE}}
\tccheck{1}
\tcstep{Send the same request for Priya Sharma}
       {HTTP 400 with `Student opted out of this meal'}
\tcstep{Send the same request for Aman Verma}
       {HTTP 403 with `You can only serve students from your own hostel'}
\tcstep{Send the same request for Dev Malhotra}
       {HTTP 422 with `Student has no wallet yet. Top up their wallet from the Wallet page first.'}
\tcstep{Send the request for Karan Patel, Lunch, day D+1}
       {HTTP 400 with `Students can only be checked in for today's meals.'}
\tcstep{Send the request for Karan Patel with the Dinner meal type, day D}
       {HTTP 400 with `Dinner check-in is only open from 19:00 to 21:30 ($\pm$60 min).'}
\tcstep{Send the request for Karan Patel without the Authorization header}
       {HTTP 401 with `Unauthorized: Missing or invalid token'}
\tccheck{2}
\end{tcsteps}

\tcend{%
  1. Rahul still has exactly one Lunch attendance record for day D and his balance is still Rs.~925.00.\newline
  2. No attendance records were created in steps 3--8, and the wallets of Priya Sharma, Aman Verma and Karan Patel are unchanged.}

% ============================================================
% 4.3 Walk-in Override
% ============================================================
\tcheader{4.3}{Walk-in Override}{Counter Check-In}
  {Test that a hostel admin can serve a walk-in student with a reason, and that Mess Staff cannot (FR-25)}

\tcbox{Pre-conditions}{%
  Postman holds tokens for Mess Staff staff.a@thapar.edu and Warden warden.a@thapar.edu, both in Hostel A.\newline
  The current time on day D is between 12:00 and 14:30.\newline
  Priya Sharma has opted out of Lunch on day D, has not been served, and has a balance of Rs.~1000.00.}

\begin{tcsteps}
\tcstep{As Mess Staff, send \texttt{POST /attendance/override} for Priya Sharma, Lunch, day D with reason `Returned early from home'}
       {HTTP 403 with `Forbidden: Insufficient privileges'}
\tcstep{As the Warden, send the same request with an empty reason}
       {HTTP 400 with a validation error for the reason field. Nothing is saved}
\tcstep{As the Warden, send the request with reason `Returned early from home'}
       {HTTP 200 with attendance result \texttt{OVERRIDE} and the reason}
\tccheck{1}
\tcstep{As the Warden, send the request from step 3 again}
       {HTTP 409 with `Student has already been served for this meal'}
\tccheck{2}
\end{tcsteps}

\tcend{%
  1. An attendance record is saved for Priya's Lunch on day D with result \texttt{OVERRIDE}, reason `Returned early from home', checked in by Neha Gupta. Priya's balance is Rs.~925.00 and a \texttt{MEAL\_DEDUCTION} of Rs.~75.00 is saved.\newline
  2. Priya still has exactly one Lunch attendance record and one deduction for day D.}

% ============================================================
% 5.1 Live Headcount Dashboard
% ============================================================
\tcheader{5.1}{Live Headcount Dashboard}{Headcount}
  {Test that expected and served counts and meal badges update live (FR-28, FR-29, FR-30)}

\tcbox{Pre-conditions}{%
  Hostel A has exactly 4 students: Rahul Kumar, Priya Sharma, Karan Patel and Dev Malhotra.\newline
  On day D, only Priya Sharma has opted out of Lunch, and nobody has been served yet.\newline
  The current time is between 12:00 and 13:30 on day D.\newline
  Browser 1 is logged in as warden.a@thapar.edu, browser 2 as staff.a@thapar.edu. Karan Patel is signed in to the mobile app.}

\begin{tcsteps}
\tcstep{In browser 1, open the Headcount page}
       {The summary shows Total Students 4, Total Booked (All Meals) 11 and Served Today 0. Lunch shows Expected 3, Actual 0, 0\% and the `Active' badge. Breakfast shows `Over' and Dinner shows `Upcoming'}
\tcstep{In browser 2, open Check-In for Lunch and click `Tap to Serve' for Rahul Kumar}
       {Rahul leaves the waiting list}
\tcstep{In browser 1, wait up to 5 seconds without reloading}
       {Lunch shows Actual 1 and 33\%. Served Today shows 1}
\tccheck{1}
\tcstep{On the mobile app, Karan Patel opts out of Dinner on day D}
       {Within 5 seconds, Dinner's Expected drops from 4 to 3 and Total Booked drops to 10}
\tcstep{After 14:30, reload the Headcount page}
       {Lunch shows the `Over' badge with Actual 1 and 33\%. Dinner shows the `Final' badge}
\tccheck{2}
\end{tcsteps}

\tcend{%
  1. The headcount for day D reports Lunch expected 3, served 1, and a total served of 1.\newline
  2. After Lunch ends, its expected and served counts are unchanged (3 and 1).}

% ============================================================
% 6.1 Wallet Top-Up
% ============================================================
\tcheader{6.1}{Wallet Top-Up}{Wallet}
  {Test that a Warden can top up a student's wallet and the student sees the new balance (FR-31, FR-32)}

\tcbox{Pre-conditions}{%
  The user is logged in to the web console as warden.a@thapar.edu.\newline
  Dev Malhotra (Hostel A) has no wallet. Aman Verma is in Hostel B.}

\begin{tcsteps}
\tcstep{Open the Wallet page}
       {The system displays the top-up form with a student search, an amount field and the semester `Autumn Semester 2026'}
\tcstep{Search `Aman'}
       {No student is found, because Aman is not in Hostel A}
\tcstep{Search `Dev', select Dev Malhotra and enter amount `0'}
       {The `Process Top-Up' button stays disabled}
\tcstep{Enter amount `1500', note `Semester advance' and click `Process Top-Up'}
       {The system confirms the top-up and clears the form}
\tccheck{1}
\tcstep{Top up Dev Malhotra again with amount `500'}
       {The system confirms the top-up}
\tccheck{2}
\tcstep{Sign in to the mobile app as Dev Malhotra and open Wallet}
       {The balance shows Rs.~2000.00 and the history shows two `Top Up' entries}
\end{tcsteps}

\tcend{%
  1. A wallet for Dev Malhotra, `Autumn Semester 2026', is saved with balance Rs.~1500.00 and one \texttt{TOPUP} transaction of Rs.~1500.00 with note `Semester advance'.\newline
  2. Dev still has one wallet, now with balance Rs.~2000.00 and two \texttt{TOPUP} transactions.}

% ============================================================
% 6.2 Semester-End Cashout
% ============================================================
\tcheader{6.2}{Semester-End Cashout}{Wallet}
  {Test that cashout is locked until the semester ends, then requested by the student and approved by the Warden (FR-33, FR-34, FR-35)}

\tcbox{Pre-conditions}{%
  Rahul Kumar's wallet for `Autumn Semester 2026' has a balance of Rs.~925.00 and he has no cashout requests.\newline
  Hostel A's fee plan `Autumn Semester 2026' has the semester end date 15/12/2026, a future date.\newline
  Rahul is signed in to the mobile app, warden.a@thapar.edu to the web console, and Postman holds tokens for both.}

\begin{tcsteps}
\tcstep{On the mobile app, open Wallet and tap the cashout button}
       {The card shows `Cashout Locked Until Semester End'. Tapping shows a `Locked' message with the semester end date}
\tcstep{In Postman, as Rahul, send \texttt{POST /wallet/cashout-request} with amount 925}
       {HTTP 403 with `Cashout is only available after your hostel's semester ends on 15 December 2026.'}
\tcstep{As the Super Admin, change the fee plan's semester end date to the day before day D. Reload Wallet on the mobile app}
       {The card shows `Cashout Available' and the button `Request Cashout of Rs.~925.00'}
\tcstep{In Postman, as Rahul, send the cashout request with amount 5000}
       {HTTP 400 with `Insufficient balance'}
\tcstep{On the mobile app, tap `Request Cashout of Rs.~925.00', then `Request'}
       {The system displays `Request Submitted'}
\tccheck{1}
\tcstep{In the web console, open the Wallet page's cashout requests and click `Approve' on Rahul's request}
       {The request leaves the pending list}
\tccheck{2}
\tcstep{In Postman, as the Warden, send \texttt{POST /wallet/cashout-requests/<id>/resolve} with status \texttt{APPROVED} for the same request}
       {HTTP 400 with `Request is already resolved'}
\end{tcsteps}

\tcend{%
  1. A cashout request for Rahul of Rs.~925.00 is saved with status \texttt{PENDING}, and his balance is still Rs.~925.00.\newline
  2. The request is \texttt{APPROVED} and resolved by Neha Gupta. Rahul's balance is Rs.~0.00 and a \texttt{CASHOUT} transaction of Rs.~925.00 is saved.}

% ============================================================
% 7.1 Hostel-Scoped Notices
% ============================================================
\tcheader{7.1}{Hostel-Scoped Notices}{Notices}
  {Test that notices reach only their own hostel and that Mess Committee members delete only their own notices (FR-39, FR-40, FR-41)}

\tcbox{Pre-conditions}{%
  Neha Gupta (Warden) and Karan Patel (Mess Committee) are in Hostel A. Rahul Kumar is a student in Hostel A and Aman Verma in Hostel B.\newline
  No notices exist.}

\begin{tcsteps}
\tcstep{In the web console as warden.a@thapar.edu, open Notices, enter title `Water supply off on Sunday' and a message, choose `Everyone (Public)' and click `Post Notice'}
       {The notice appears at the top of the list with Neha Gupta's name and a Warden badge}
\tccheck{1}
\tcstep{On the mobile app as Rahul Kumar, open Notices}
       {The notice `Water supply off on Sunday' is shown}
\tcstep{On the mobile app as Aman Verma, open Notices}
       {The notice is not shown}
\tcstep{In the web console as karan@thapar.edu, post the notice `Menu change tomorrow'}
       {Karan's notice shows a delete icon. Neha's notice shows no delete icon}
\tcstep{Click the delete icon on Karan's notice}
       {Karan's notice is removed from the list}
\tccheck{2}
\end{tcsteps}

\tcend{%
  1. The notice `Water supply off on Sunday' is saved for Hostel A with no target role, posted by Neha Gupta.\newline
  2. Karan's notice is deleted and Neha's notice still exists.}

\end{document}
